\section{The generalised winding number}\label{sec:winding}

The classical winding number is defined for curves that avoid the
reference point. To accommodate curves that pass through that point,
Hungerbühler--Wasem use a Cauchy principal value of the Cauchy
kernel.

\begin{definition}[Generalised winding number]
  \label{def:generalized-winding-number}
  \lean{generalizedWindingNumber}
  \uses{def:has-cauchy-pv-on}
  \leanok
  Let $\gamma : [0,1] \to \mathbb{C}$ be a piecewise-$C^{1}$ path
  from $x$ to $y$, and let $z_0 \in \mathbb{C}$. The
  \emph{generalised winding number} of $\gamma$ around $z_0$ is
  defined by
  \[
    \operatorname{gWN}(\gamma, z_0)
    \;=\; \frac{1}{2 \pi i}
    \operatorname{PV}\!\!\int_{\gamma}
      \frac{\mathrm{d}z}{z - z_0},
  \]
  where the principal value is the single-point CPV at $z_0$ (a
  specialisation of \cref{def:has-cauchy-pv-on} with
  $S = \{z_0\}$ and $f(z) = (z-z_0)^{-1}$). When the principal value
  does not exist, $\operatorname{gWN}(\gamma, z_0)$ is set to a
  default value.

  When $\gamma$ avoids $z_0$, $\operatorname{gWN}(\gamma, z_0)$
  coincides with the classical winding number, an integer. When
  $\gamma$ crosses $z_0$ transversally once, the principal value
  exists with $\operatorname{gWN}(\gamma, z_0) \in \tfrac{1}{2}
  + \mathbb{Z}$. More generally,
  $\operatorname{gWN}(\gamma, z_0)$ may take any half-integer or
  rational value, depending on how often and at what angles $\gamma$
  meets $z_0$.
\end{definition}
